% =====================================================================
% ARTIGO — PROJETO DE ENGENHARIA DE SOFTWARE (Prof. Sérgio Santana)
% Tema 6 · Teste de usabilidade em aplicativos (estudo de caso)
% Caso 15 · Reserva de laboratórios da faculdade
% ---------------------------------------------------------------------
% Base: templates/modelo-tcc-sbc.tex (modelo SBC fornecido pelo docente).
% Ajustes exigidos pelo PDF da disciplina (seção 6):
%   - seções: título, resumo, introdução, fundamentação teórica,
%     desenvolvimento, considerações e referências;
%   - Times New Roman 12, espaçamento 1,5;
%   - de 8 a 10 páginas, mínimo de 5 referências (2 artigos científicos);
%   - entrega em PDF em 30/11.
% Citações no formato autor-data (ABNT).
% Trechos ainda não escritos aparecem marcados com \pendente{...}.
% =====================================================================

\documentclass[a4paper,12pt]{article}

% ---------------------------------------------------------------------
% PACOTES (os mesmos do modelo SBC + setspace, xcolor e tabularx)
% ---------------------------------------------------------------------
\usepackage[utf8]{inputenc}
\usepackage[T1]{fontenc}
\usepackage[brazilian]{babel}
\usepackage{times}
\usepackage[a4paper,top=3.5cm,left=3cm,right=3cm,bottom=2.5cm]{geometry}
\usepackage{graphicx}
\usepackage{url}
\usepackage{indentfirst}
\usepackage{titlesec}
\usepackage[bf,sf,footnotesize,indent]{caption}
\usepackage{setspace}
\usepackage{xcolor}
\usepackage{tabularx}

% ---------------------------------------------------------------------
% ESPAÇAMENTO: 1,5 entre linhas (regra da disciplina)
% ---------------------------------------------------------------------
\onehalfspacing
\setlength{\parindent}{1.27cm}
\setlength{\parskip}{6pt}

% ---------------------------------------------------------------------
% SEÇÕES E LEGENDAS (iguais ao modelo SBC)
% ---------------------------------------------------------------------
\titleformat{\section}{\fontsize{13}{15}\bfseries}{\thesection.}{1em}{}
\titleformat{\subsection}{\fontsize{12}{14}\bfseries}{\thesubsection.}{1em}{}
\titleformat{\subsubsection}{\fontsize{12}{14}\bfseries\itshape}{\thesubsubsection.}{1em}{}
\titlespacing*{\section}{0pt}{12pt}{6pt}
\titlespacing*{\subsection}{0pt}{10pt}{4pt}

\renewcommand{\captionfont}{\sffamily\footnotesize\bfseries}
\setlength{\captionmargin}{0.8cm}
\captionsetup{labelsep=period}

% ---------------------------------------------------------------------
% RESUMO / ABSTRACT (iguais ao modelo SBC)
% ---------------------------------------------------------------------
\newenvironment{resumo}{%
  \list{}{\advance\topsep by6pt\relax
    \leftmargin=0.8cm
    \labelwidth=0pt
    \listparindent=0pt
    \itemindent\listparindent
    \rightmargin\leftmargin}\item[\hskip\labelsep\bfseries\itshape Resumo.]\itshape}%
  {\endlist}

\renewenvironment{abstract}{%
  \list{}{\advance\topsep by6pt\relax
    \leftmargin=0.8cm
    \labelwidth=0pt
    \listparindent=0pt
    \itemindent\listparindent
    \rightmargin\leftmargin}\item[\hskip\labelsep\bfseries\itshape Abstract.]\itshape}%
  {\endlist}

% Marca de trecho pendente (remover todas antes da entrega final)
\newcommand{\pendente}[1]{{\color{gray}\upshape [PENDENTE: #1]}}

% Referências ABNT: sem número, alinhadas à esquerda
\newenvironment{referencias}{%
  \section*{Referências}
  \begin{singlespace}
  \setlength{\parindent}{0pt}\setlength{\parskip}{10pt}\raggedright}%
  {\end{singlespace}}

\begin{document}
\pagestyle{empty}

% =====================================================================
% TÍTULO, AUTOR, RESUMO E ABSTRACT
% =====================================================================
\begin{center}
  \begin{singlespace}
  {\fontsize{16}{20}\selectfont\bfseries Teste de Usabilidade com Cinco Usuários: Estudo de Caso de um Sistema de Reserva de Laboratórios}\\[4pt]
  {\fontsize{14}{18}\selectfont\bfseries Usability Testing with Five Users: A Case Study of a Laboratory Reservation System}\\[12pt]

  {\normalsize\bfseries Wilia Francisco da Silva$^{1}$}\\[6pt]
  {\footnotesize\itshape Disciplina: Laboratório de Engenharia de Software -- Prof. Sérgio Santana}\\[4pt]

  {\normalsize
    $^{1}$Curso de Ciência da Computação -- Universidade Tiradentes (UNIT) -- Aracaju/SE\\{}
    \footnotesize\texttt{wilia.francisco@souunit.com.br}
  }
  \end{singlespace}
\end{center}
\vspace{12pt}

\begin{resumo}
\pendente{escrever ao final, em um único parágrafo, com objetivo,
método, principais resultados e conclusão.}
\end{resumo}
\noindent\textbf{Palavras-chave:} usabilidade, teste com usuários, System Usability Scale, avaliação heurística.

\vspace{6pt}

\begin{abstract}
\pendente{versão em inglês do resumo.}
\end{abstract}
\noindent\textbf{Keywords:} usability, user testing, System Usability Scale, heuristic evaluation.

\vspace{18pt}

% =====================================================================
% 1. INTRODUÇÃO (entregue na Aula 3)
% =====================================================================
\section{Introdução}

Um sistema pode cumprir todas as funções previstas e, ainda assim,
fracassar quando as pessoas que o usam não entendem o que fazer,
cometem erros com frequência ou desistem da tarefa. Por isso, além de
verificar se o software funciona, é preciso verificar se ele é fácil
de usar. Segundo Barbosa e Silva (2010), a usabilidade está relacionada
à facilidade de aprendizado e de uso de um sistema e à satisfação de
quem o utiliza, e deve ser avaliada com métodos próprios ao longo do
desenvolvimento.

Esse cuidado vale também para sistemas pequenos, criados para resolver
problemas do dia a dia. É o caso do sistema de reserva de laboratórios
de uma faculdade: a coordenação recebia os pedidos por e-mail e era
comum dois professores chegarem ao mesmo laboratório no mesmo horário.
A versão 1.0 do sistema foi criada para mostrar as reservas e impedir
esses choques. No entanto, se os professores tiverem dificuldade para
reservar, a tendência é que voltem a usar o e-mail, e o problema
original continua.

Avaliar a usabilidade não exige, necessariamente, muitos participantes.
O modelo proposto por Nielsen e Landauer (1993) mostra que a quantidade
de problemas encontrados cresce rapidamente com os primeiros usuários
testados e que um grupo pequeno, por volta de cinco pessoas, já revela
a maior parte deles. A observação dos participantes pode ser combinada
com um questionário de satisfação, como a escala SUS (BROOKE, 1996), e
com a classificação dos problemas pelas heurísticas de usabilidade
(NIELSEN; MOLICH, 1990).

Diante disso, este artigo busca responder à pergunta: \emph{como testar
se um aplicativo é fácil de usar?} Para isso, apresenta um estudo de
caso em que a versão 1.0 do sistema de reserva de laboratórios foi
testada por cinco colegas de curso. Os objetivos específicos são:
(i) planejar tarefas que representem o uso real do sistema;
(ii) medir a taxa de sucesso e o tempo de cada tarefa;
(iii) aplicar o questionário SUS; e
(iv) classificar os problemas encontrados e propor melhorias para a
versão 2.0.

O restante do texto está organizado da seguinte forma: a Seção 2
apresenta a fundamentação teórica sobre usabilidade e testes com
usuários; a Seção 3 descreve o desenvolvimento do estudo e analisa os
resultados; e a Seção 4 traz as considerações finais.

% =====================================================================
% 2. FUNDAMENTAÇÃO TEÓRICA (Aula 4 · 23/11)
% =====================================================================
\section{Fundamentação Teórica}

\subsection{Usabilidade}

Barbosa e Silva (2010) apresentam a usabilidade como um dos principais
critérios de qualidade de uso de um sistema interativo. Com base na
norma ISO 9241-11, os autores a descrevem como o grau em que um produto
pode ser usado por usuários específicos para alcançar objetivos
específicos com eficácia, eficiência e satisfação, em um contexto de uso
determinado. A eficácia indica se o usuário consegue concluir a tarefa;
a eficiência, quanto tempo e esforço ele gasta para isso; e a
satisfação, como ele se sente ao usar o sistema.

Essas três dimensões mostram que um sistema não é fácil de usar apenas
por funcionar: ele
precisa ser compreendido logo no primeiro uso, evitar que o usuário
cometa erros e permitir que ele se recupere quando os erros acontecem.
No sistema de reserva de laboratórios, por exemplo, uma regra correta,
como o limite de 1 a 4 horas por reserva, pode gerar frustração se o
professor só descobrir esse limite depois de preencher todo o
formulário.

\subsection{Teste de usabilidade com usuários}

Barbosa e Silva (2010) distinguem os métodos de avaliação feitos por
especialistas, que inspecionam a interface, dos métodos de observação,
em que pessoas representativas do público-alvo usam o sistema. No teste
de usabilidade, o avaliador prepara tarefas realistas, observa os
participantes enquanto eles as executam e registra dados como o sucesso
na tarefa, o tempo gasto e os erros cometidos. É comum pedir que o
participante pense em voz alta, explicando o que está tentando fazer, o
que ajuda a entender a causa das dificuldades, e não apenas o seu
efeito.

A quantidade de participantes é uma dúvida frequente. Nielsen e
Landauer (1993) propuseram um modelo matemático em que a proporção de
problemas encontrados depois de $i$ participantes é
$1-(1-\lambda)^{i}$, sendo $\lambda$ a probabilidade de um único
participante revelar um problema qualquer. Nos projetos analisados
pelos autores, $\lambda$ ficou, em média, em torno de 31\%. Com esse
valor, cinco participantes revelam cerca de 85\% dos problemas, e cada
participante a mais traz cada vez menos descobertas novas. Os autores
também mostram, numa análise de custo e benefício, que o melhor retorno
está em testes com poucos participantes, repetidos depois de cada
rodada de correções. O resultado, porém, depende de $\lambda$, que muda
de um sistema para outro; os cinco participantes são uma referência
prática, e não uma garantia.

\subsection{Questionário SUS}

A \emph{System Usability Scale} (SUS) foi proposta por Brooke (1996)
como uma forma rápida e barata de medir a percepção de usabilidade de um
sistema. O questionário tem dez afirmações, respondidas numa escala de 1
(discordo totalmente) a 5 (concordo totalmente), alternando afirmações
positivas (itens ímpares) e negativas (itens pares). Para calcular a
nota, subtrai-se 1 da resposta de cada item ímpar e subtrai-se de 5 a
resposta de cada item par; a soma desses valores é multiplicada por
2,5, o que gera uma nota de 0 a 100 (BROOKE, 1996). Essa nota não é uma
porcentagem: ela só tem sentido quando comparada a uma referência.

Essa referência foi construída por Bangor, Kortum e Miller (2008), que
analisaram mais de dois mil questionários SUS aplicados em cerca de
duzentos testes de usabilidade, ao longo de dez anos. Os autores
confirmaram a alta confiabilidade do questionário e encontraram uma nota
média próxima de 70. Com base nisso, propuseram faixas de
aceitabilidade: notas abaixo de 50 indicam usabilidade inaceitável;
entre 50 e 70, marginal; e acima de 70, aceitável (BANGOR; KORTUM;
MILLER, 2008). Neste estudo, essas faixas são usadas para interpretar a
nota obtida pelo sistema.

\subsection{Heurísticas de usabilidade}

Nielsen e Molich (1990) propuseram a avaliação heurística, em que
avaliadores examinam a interface com base em um pequeno conjunto de
princípios gerais, chamados heurísticas. As nove heurísticas originais
são: diálogo simples e natural; falar a língua do usuário; minimizar a
carga de memória do usuário; ser consistente; fornecer \emph{feedback};
oferecer saídas claramente marcadas; oferecer atalhos; fornecer boas
mensagens de erro; e prevenir erros. Os autores observaram que cada
avaliador, sozinho, encontrou apenas de 20\% a 51\% dos problemas, e
que juntar os resultados de vários avaliadores aumenta bastante a
quantidade de problemas encontrados (NIELSEN; MOLICH, 1990).

Neste trabalho, as heurísticas não são usadas como um método de
avaliação separado, mas como vocabulário para classificar os problemas
observados no teste com usuários. Cada problema também é classificado
pelo impacto que teve na tarefa observada: \emph{alto}, quando impediu o
participante de concluí-la; \emph{médio}, quando a atrasou ou exigiu uma
nova tentativa; e \emph{baixo}, quando causou apenas dúvida ou
incômodo. Essa classificação ajuda a decidir quais correções devem
entrar primeiro na versão 2.0.

\subsection{Síntese}

Combinadas, essas referências orientam o método deste estudo: as tarefas
observadas medem a eficácia (taxa de sucesso) e a eficiência (tempo) do
uso; o questionário SUS mede a satisfação; e as heurísticas, com a
classificação por impacto, organizam os problemas encontrados numa
lista de melhorias para a versão 2.0.

% =====================================================================
% 3. DESENVOLVIMENTO (Aula 5 · entrega até 22h)
% =====================================================================
\section{Desenvolvimento}

\subsection{O sistema avaliado}
\pendente{contexto do Caso 15, regras principais e telas da versão 1.0
(pode usar a figura da arquitetura em 3 camadas da Aula 3).}

\subsection{Planejamento do teste}
\pendente{perfil dos 5 participantes, local e equipamento, tarefas,
roteiro, termo de consentimento e métricas coletadas.}

\subsection{Resultados}
\pendente{tabela com sucesso e tempo por tarefa; nota SUS de cada
participante e média; lista de problemas encontrados.}

\subsection{Análise e propostas para a versão 2.0}
\pendente{classificação dos problemas por heurística e por impacto
(alto, médio ou baixo) e a melhoria proposta para cada um.}

% =====================================================================
% 4. CONSIDERAÇÕES FINAIS (Aula 6 · 30/11)
% =====================================================================
\section{Considerações Finais}
\pendente{resposta à pergunta do artigo, principais achados,
limitações (participantes eram colegas no papel de professores) e
trabalhos futuros.}

% =====================================================================
% DECLARAÇÃO DE USO DE IA (obrigatória no modelo SBC)
% =====================================================================
\section*{Declaração de Uso de Inteligência Artificial}
\pendente{informar a ferramenta de IA usada, em quais etapas do
trabalho e com qual finalidade.}

% =====================================================================
% REFERÊNCIAS (ABNT, ordem alfabética)
% =====================================================================
\begin{referencias}

BANGOR, A.; KORTUM, P. T.; MILLER, J. T. An empirical evaluation of the
System Usability Scale. \emph{International Journal of Human-Computer
Interaction}, v. 24, n. 6, p. 574--594, 2008. DOI: 10.1080/10447310802205776.

BARBOSA, S. D. J.; SILVA, B. S. \emph{Interação humano-computador}. Rio
de Janeiro: Elsevier, 2010.

BROOKE, J. SUS: a ``quick and dirty'' usability scale. In: JORDAN, P. W.
et al. (Ed.). \emph{Usability Evaluation in Industry}. London: Taylor \&
Francis, 1996. p. 189--194.

NIELSEN, J.; LANDAUER, T. K. A mathematical model of the finding of
usability problems. In: \emph{Proceedings of the INTERCHI'93 Conference
on Human Factors in Computing Systems}. Amsterdam: ACM, 1993.
p. 206--213. DOI: 10.1145/169059.169166.

NIELSEN, J.; MOLICH, R. Heuristic evaluation of user interfaces. In:
\emph{Proceedings of the SIGCHI Conference on Human Factors in Computing
Systems (CHI'90)}. Seattle: ACM, 1990. p. 249--256. DOI: 10.1145/97243.97281.

\end{referencias}

\end{document}
